\chapter{异常检测}
\chaptergoal{区分异常分数、算法候选和真实异常身份；理解 Isolation Forest 的随机分裂思想；观察阈值如何影响候选数量；使用带已知人工异常的受控实验检验程序；在天文与物理数据中建立有记录的候选复核流程。}

在科学研究中，异常对象往往引人注意：可能是有价值的特殊天体，可能是罕见实验状态，也可能只是测量、匹配或处理错误。异常检测可以帮助缩小需要人工核查的范围，但它不能替代科学判别。教材必须把“算法认为不寻常”和“已经证实的新现象”分开。
\section{异常检测的任务与数据条件}

\textbf{异常检测（anomaly detection）}是在指定数据范围和参照规则下，找出不符合常见模式的记录。它首先回答“值得优先检查哪些记录”，不是自动给出异常的物理原因。\textbf{离群点}通常强调记录偏离当前数据主体；\textbf{新颖性检测}则常先用相对正常的参考数据学习，再判断后来出现的记录是否偏离参考模式。具体接口需要与任务设定对应。

例如多数记录的温度和压强都在某个范围内，一条温度极高的记录可能是一种单变量异常；另一条温度和压强各自都不极端，但组合从未出现过，也可能是多变量异常。二者都可能来自真实状态，也可能来自单位混乱或测量错误。在调用检测器之前，必须先检查字段含义。

异常通常指与多数记录或模型预期显著不同的样本。它可以是单变量的极端数值，也可以是多变量空间中非常罕见的特征组合。不同检测方法隐含不同正常性定义：距离方法关心远离常见记录的程度，密度方法关心局部密度，而残差方法关心偏离已有预测模型的程度。

本章主要讨论无监督异常检测。若已经有大量可靠的“正常/异常”标签，也可以训练监督分类器；若训练时只有正常样本，则又属于新颖性检测的一种场景。不同任务对应不同的训练与评价方式，不应统一称为“无监督”而忽视标签是否被使用。

\subsection{异常是相对于参照与表示而言}
一个数值可能相对整体分布很罕见，却在某一实验工况下完全正常。温度很高时传感器输出较大，若模型忽略温度，就可能把正常响应误报为异常。因此异常检测首先要说明参照总体、输入表示和应用目标，而不是先选择某个分数函数。

可区分单点的极端值、给定条件下的不寻常记录，以及只有一组记录共同出现才显著的异常。表格算法主要处理当前所给特征的模式，若关心时序中的突变，却没有提供时间或变化特征，算法未必能感知所需结构。

\textbf{离群点检测}常在可能混有异常的现有数据中寻找不一致记录；\textbf{新颖性检测}先利用相对可信的参照数据建立模型，再判断后来的记录是否偏离。两者都可不使用异常类别标签，但训练与评价的数据角色不同。某些软件接口对二者有明确区别，尤其是LOF，不能直接把训练上的分数当作新样本预测规则。

\subsection{异常检测与稀有类别分类不完全相同}
若有充分、可靠的异常标签，并且未来异常类型与训练相似，可以建立监督分类模型。无监督异常方法则主要依据分布结构，不保证找到某种指定物理类别。一个科学上极重要的稀有天体可能在当前测量特征上并不离群；一个非常离群的记录也可能只是单位错误。

因此要分开“发现值得检查的记录”与“确认一种科学身份”。本章训练前者，确认需要原始数据、仪器信息、独立观测或其他证据。

\section{异常分数、阈值与候选数量}

\textbf{异常分数}是按算法规则计算的排序量；\textbf{阈值}是把这个连续量转为“入选或不入选”的分界。异常分数通常不是“属于异常的概率”，除非另有明确概率模型或校准证据。

设五条人工记录的分数依次为0.10、0.20、0.35、0.80、0.90，约定越大越异常。用0.75作阈值会选出最后两条；改用0.30会选出后三条。假设在这个受控例子中已知最后两条才是注入对象，前一种阈值有两次检出、零次误报，后一种有两次检出、一次误报。实际数据没有这样完整的已知身份时，就只能报告候选数量及已经复核的证据，不能凭分数自行产生真值。

检测器常给每条记录一个连续分数，再用阈值挑选候选。分数的正负方向由具体接口定义，不同算法之间不能不经变换直接比较数值。以 sklearn 的 IsolationForest 为例，\texttt{score\_samples()} 越低通常表示越异常；\texttt{decision\_function()} 在默认偏移调整后也通过较低值表示更偏异常的对象。使用前须检查对应 API 文档和实际返回值。

阈值越宽松，通常候选越多，但候选是否真的值得科学关注取决于外部验证。若设置 \texttt{contamination=0.05}，它主要影响阈值偏移和预期异常比例的界定，不代表已经知道数据中恰好有5\%真实异常。课堂可以比较多个阈值，绘制候选数量的变化。

\subsection{排序、决策与预算是三个对象}
异常分数给出相对不寻常程度，阈值把连续分数转成候选与非候选，检查预算决定能复核多少记录。分数高低方向由实现规定：有些接口高分更正常，有些高分更异常。本书实验统一记录高值表示更异常的变量，并在调用时明确取负号或变换，避免排序反了仍得到一张看似合理的图。

给定每周只能人工检查20条记录，可以按分数取前20，而不必声称这些记录超过某个普适概率阈值。此时评价可以关注前20中的已确认异常数或已知注入的召回。预算型选择与固定误报率选择服务于不同目标，应提前说明。

\worked{很小的误报率仍可能产生很多候选}
假设筛查10万条记录，其中真实关注异常100条。若召回率80\%、正常记录误报率1\%，则真异常检出80条，正常误报约999条，候选共1079条，精确率约7.4\%。这一人工算例说明罕见性会影响候选纯度，不能只报告“误报率仅1\%”就宣称很可靠。

改变阈值通常在检出与检查成本之间取舍。异常分数不一定是概率，阈值0.5也不自动表示“异常概率超过一半”。如果使用参考数据某个分位数定阈值，这只是依据该参照分布的操作规则；未来分布变化会改变误报率。

\subsection{校准阈值需要考虑训练使用范围}
用同一批记录训练异常模型再取其分位数，得到的是训练参照分数的阈值，并不能直接保证独立新记录有相同超阈比例。更严格做法是在不参与模型拟合的正常校准记录上选择阈值，再对新记录检查。若正常参照本身含异常，阈值也会受到污染。报告须明确阈值来自训练、校准还是某个科学规定。

\section{Isolation Forest 的直观原理}

\textbf{孤立森林（Isolation Forest）}把“一个点是否容易被随机切分单独隔离”作为异常倾向的依据。\textbf{路径长度}就是一条记录从根节点到单独隔离或停止节点经历的切分次数。

用人工一维点$0,0.1,0.2,5$说明：若第一次随机阈值落在2，点5立即独自位于右侧，只需一次分裂；左侧三点还需更多分裂才能区分。一次随机阈值并不能证明谁异常，所以森林汇总多棵随机树的隔离路径，并按子样本规模作相应处理。这个例子说明短路径的直觉，不是用单次阈值代替真实算法的分数公式。\textsuperscript{[23]}

Isolation Forest 通过随机选择特征及分裂位置，把样本逐步隔离出来。较容易用少数随机分裂孤立的样本，往往被赋予更强的异常倾向。它与逐点计算所有距离的办法不同，对一些规模较大的数据也具有使用价值，但其异常性仍由当前输入空间定义。

\fig{ch15_isolation.pdf}{人工生成二维数据及 Isolation Forest 的异常分数示意。人工注入的特殊点以独立标记显示，图不代表实际天体或实验测量。}

在极端维度、特征尺度差异或存在多个相互分离正常群体的情况下，模型可能把边缘正常对象视作异常。因此需要通过已知控制实验和专业信息解释分数，而不是只选择最极端的几个点就宣布发现。

\subsection{从随机切分到平均隔离路径}
隔离森林不是通过监督标签寻找“最纯”的节点，而是在子样本上随机选择特征与切分位置，逐步把记录隔离。处于稀疏或极端位置的点，往往较少切分就能被单独分出；密集区域内的记录通常需要更多切分。每棵树给出一条隔离路径长度，多棵树取平均以减少随机切分的偶然性。

以一维人工点 $0,0.1,10$ 为例，在最小到最大值之间均匀取第一刀。若阈值落在0.1与10之间，点10立即单独成为一侧，而两个近邻点仍在一起。这个区间占总范围的大部分，所以极端点更容易被快速隔离。若只是离整体中心远、但自己构成一个密群，路径行为会不同，不能把隔离分数等同于到均值的距离。

\subsection{路径归一化为什么不是概率模型}
常见隔离评分把平均路径长度 $E[h(x)]$与样本规模为 $\psi$ 的参考路径长度 $c(\psi)$比较，例如
\[
s(x)=2^{-E[h(x)]/c(\psi)},\qquad
c(\psi)=2H_{\psi-1}-2(\psi-1)/\psi,
\]
其中 $H_m=\sum_{j=1}^{m}1/j$。这是基于随机树路径的归一化形式；实现还会对未完全隔离的叶子作路径修正。路径较短对应较大 $s$，但 $s$不是根据异常标签估计的后验概率。

对 $\psi=3$，$c=5/3$。若某点平均路径为1，分数约 $2^{-0.6}=0.6598$；若为2，约 $2^{-1.2}=0.4353$。两者可以比较隔离程度，却不能说前者“有65.98\%概率异常”。实际sklearn的 \texttt{score\_samples}采用与这个高异常方向相反的符号约定，本书显式使用其负值作高异常分数。[23]

\subsection{污染比例与模型配置的作用}
\texttt{contamination}通常影响把分数转成异常标签的阈值设置，不是告诉算法已知哪些行是异常。设置0.05不能让模型自动识别真实5\%异常，也不是从数据证明异常比例恰为5\%。树数、子样本数和可用特征影响隔离结构，阈值与结构应分开研究。

若无关特征很多，随机切分常发生在无关方向，真正异常结构可能被稀释。若一个标识列近乎连续而无科学意义，它也可能影响隔离。使用无监督方法并不能免除字段角色检查、单位核对和相关对象处理。

\section{LOF 与局部密度的选读比较}

\textbf{局部密度}描述一个点附近的记录有多集中。假设数据的一部分本来很密，另一部分整体较稀，直接比较每点到最近邻的距离可能把整个稀疏群体都判为异常。LOF转而把某点的局部密度与其邻居的密度相比较，关注它是否相对自己的邻域特别稀疏，而不仅仅看它离全体均值多远。该机制仍取决于选择了多少邻居以及用什么距离。

局部异常因子（local outlier factor, LOF）通过局部邻域密度比较评估对象异常程度。它可能识别“总体不算很远，但相对附近样本很稀疏”的点。邻居数与距离定义会影响判断，尺度与高维问题仍然存在。sklearn 中 LOF 的普通离群检测模式与新颖性预测模式在接口用途上有区别，课堂必须根据所用模式调用方法。

本章以 Isolation Forest 为主线，LOF 作为概念比较与提高题；两者的分数不同，不应直接用数值大小横向排名。

\subsection{局部参照与全局参照为何不同}
局部离群因子LOF将一个点附近的密度与其邻居的密度比较。一个位于整体稀疏但内部正常的群体中的点，不一定异常；相反，一个嵌在密群边缘的小空隙中、比邻居明显稀疏的点可能具有较大LOF。邻居数量规定比较尺度，过小对噪声敏感，过大可能把不同群体混在一起。

令 $k$距离为某点到第 $k$近邻的距离。对点 $p$与邻居 $o$，\textbf{可达距离}定义为
\[
\mathrm{reach}_k(p,o)=\max\{d_k(o),d(p,o)\}.
\]
它使用邻居 $o$ 的 $k$距离作下限，减少极近点对造成的不稳定。局部可达密度是平均可达距离的倒数，LOF则为邻居密度相对本点密度的平均比例。接近1表示与邻居相当，明显大于1表示本点相对稀疏；有限样本与并列距离需要按具体定义处理。

\worked{四个一维点的一次LOF手算}
取 $0,1,2,10$，$k=2$。点10的两个邻居是2和1，距离8和9；邻居2的 $k$距离为2，邻居1的为1，因此点10的可达距离仍为8和9，局部密度为 $2/17$。

点2的邻居为1和0，可达距离分别为1和2，密度为 $2/3$；点1的邻居为0和2，两者的 $k$距离均为2，故可达距离均为2，密度为 $1/2$。于是
\[
\mathrm{LOF}(10)=\frac{(2/3+1/2)/2}{2/17}=\frac{119}{24}\approx4.9583.
\]
这个点的密度显著低于两个邻居。计算必须使用可达距离，若直接以最近距离倒数代替，会改变算法。

\subsection{对新样本的接口限制}
普通LOF用于现有样本离群分析时，训练记录参与彼此的邻域定义。需要判断后来新对象时，应明确采用新颖性模式，并只对未参与拟合的记录调用相应预测接口。用同一组记录同时作为训练邻域与新查询，会改变排除自身等细节，不能假设结果完全等价。[23]

\fig{v3_ch15_local_density.pdf}{人工数据中局部密度结构的示意。不同密度群体需要局部比较，不能仅依据到全体均值的距离判断异常。}

\section{受控异常实验}
\begin{lstlisting}[caption={人工生成二维数据上的异常检测（仅用于方法实验）}]
import numpy as np
from sklearn.ensemble import IsolationForest

rng = np.random.default_rng(42)
normal = rng.normal(size=(240, 2))
injected = rng.uniform(-6, 6, size=(15, 2))
X = np.vstack([normal, injected])
truth = np.r_[np.zeros(240, dtype=int), np.ones(15, dtype=int)]
model = IsolationForest(n_estimators=150,
                        contamination="auto", random_state=42)
score = -model.fit(X).score_samples(X)  # higher = more anomalous
for threshold in np.quantile(score, [0.90, 0.95, 0.98]):
    candidates = score >= threshold
    print(threshold, candidates.sum(),
          ((candidates == 1) & (truth == 1)).sum())
\end{lstlisting}

\texttt{truth}记录的是生成来源：0表示正态主体样本，1表示后来加入的均匀分布样本。注意均匀注入的范围覆盖主体区域，其中一些注入点可能并不远离主体；正态样本本身也可能出现尾部点。因此这个对照严格说是在检验“识别不同生成来源”的能力，不能把每一个来源为1的点都视为已经证明的几何异常。因为来源已知，可以统计被选中的注入数量以及主体记录的误选数量。但真实科学数据通常没有完整的异常真值表：即使一个候选非常偏离，也可能是数据处理错误；即使没有被算法标出，也可能是有趣而不极端的物理现象。

\fig{ch15_threshold.pdf}{在受控异常实验中，改变分数阈值对应的候选数量及人工注入对象的检出数；结果来自固定随机状态。}

\subsection{把生成、训练、阈值与评价分开}
本轮实验生成400条二维标准正态参考记录，用于训练隔离森林。另生成200条同分布正常查询和20条位于较远矩形区域的人工异常，形成220条查询。随机状态固定为150；这些不是物理测量，也不是天文发现。异常标签仅用于查询集的检验，不进入无监督训练。

定义高异常分数为 $-\texttt{score\_samples}$，使用参考训练分数的90\%、95\%、99\%分位数定三个候选阈值。这里阈值来自训练参照，仅作为机制实验，不宣称对未来记录严格控制某个误报率。结果如下：
\begin{center}\small
\begin{tabular}{rrrrr}\toprule
参考分位数&阈值&候选数&异常召回率&候选精确率\\\midrule
0.90&0.5503&40&1.0000&0.5000\\
0.95&0.5920&29&1.0000&0.6897\\
0.99&0.6685&20&1.0000&1.0000\\\bottomrule
\end{tabular}\end{center}
\fig{v3_ch15_scores.pdf}{独立生成查询记录的异常分数与注入身份。异常被故意放在较远区域，这是一项容易分离的受控检验，不代表真实异常任务都能达到相同指标。}

三个阈值都找到20条人工异常，但较低阈值分别多报20和9条正常记录。最高阈值在本查询集中恰好没有误报。这个结论只描述有限的220条查询，不能由零误报估计总体误报率必为0。

\subsection{一个有效实验也应说明它没有检验什么}
本例异常与正常分得较远，所以平均精确率达到1并不令人意外。更严格的提高实验可以把异常逐渐移近正常分布、仅改变相关结构、加入不相关维度，或让参照混入少量异常。每次只改变一个机制，记录召回、误报与分数分布。程序应保留生成参数和注入行号，而不是为得到理想结果反复选择随机种子。

对真实数据，通常不知道完整异常身份，不能直接制作同样的混淆矩阵。可以报告人工复核覆盖率、已确认候选比例以及核验规则，同时明确未复核对象仍未知。受控注入用于验证某种异常模式的检出能力，不等于测量全部未知异常的完备性。

\section{专业候选的复核与记录}
天文案例可选周期--星等关系中残差较大的变量星，或者 Gaia 候选结构附近特征特殊的对象。首先核对源身份和质量标记，再检查相关测量、相邻对象、重复记录和原始观测信息。物理实验则检查校准状态、工况记录、仪器饱和和重复测量。

输出候选清单应包含唯一身份、异常分数、阈值、输入特征、质量与复核结论。未核实的样本应标为“待检查候选”。对稀有对象的实际检出率，若缺乏可靠标签不能虚构一个数值；可以报告受控注入实验的检测能力及其适用限制。

\subsection{从异常分数到证据表}
每个候选保留对象标识、原始测量、输入变换、分数、阈值、质量标志和复核结论。复核至少区分：数据格式或单位错误；仪器和处理缺陷；已知罕见但合理对象；尚不能解释的候选。未知不能直接填成新发现，未确认也不能自动算作误报。

天文光谱残差异常时，应查看原光谱、掩码、信噪比、速度对齐和邻近谱线；物理实验异常时，应核对工况、传感器量程、批次和原始日志。一个模型分数不能替代这些证据。对成本较高的跟进，可以按分数、质量和科学价值共同安排优先级，但要明确这已经是另一个决策规则。

\subsection{复核反馈可以形成下一轮监督问题}
人工核验积累可靠标签后，可以研究针对某种已知异常的监督分类器。此时应重新定义正负类、抽样和评价，不能直接把只核验高分候选得到的标签当作无偏总体样本。因为低分区域没有同等核验机会，反馈数据本身有选择偏差。

异常检测课程的成果不是一张颜色很醒目的散点图，而是一条可解释、可追溯的候选处理流程：参照如何建立，分数怎样计算，阈值服务什么预算，哪些对象被检查，最后的结论得到什么证据支持。

\chapterpractice{在受控生成数据中训练 Isolation Forest，比较不同阈值下的候选数与已知注入的检出；选一份专业数据建立复核记录表，说明哪些异常仍可能来自数据质量或测量条件。}
\chapterexercise
\begin{enumerate}
\item 在已约定越高越异常的分数中，数值很高能否证明某个对象是新天体类型？
\item 如果人为规定选出5\%的记录，是否意味着真实异常率就是5\%？
\item 比较整体距离离群和局部密度异常的差别。
\item 为什么必须将人工注入异常与原始真实数据明确分开保存？
\end{enumerate}

\section*{深化练习与方法检验}
本组题强调解释、计算与验证；D编号在本章内独立计数。计算题应写出所用定义和中间步骤，实验题保留实际输出，包括不支持预期的结果。
\begin{enumerate}[label=D\arabic*.]
\item 参考正常样本的分数95\%分位数被用作阈值，是否能保证任意未来数据上只有5\%误报？区分经验分位数、总体误报率和数据分布变化。
\item 在本章受控实验中，提高阈值后精确率可达1、召回率仍为1。这能否证明检测器适合真实天文异常？怎样使注入检验更有诊断价值？
\item 对于一维点0、1、2、10，取2近邻并按正文LOF定义计算点10的局部离群因子。为什么不能只用“到最近点的距离”代替局部密度比？
\item 候选清单只允许人工检查20个对象。怎样报告这种预算约束下的结果，为什么不能把“前20名”叫作全部真实异常？
\end{enumerate}
